\documentclass[a4paper,12pt]{article}

\usepackage[T2A]{fontenc}
\usepackage[utf8]{inputenc}
\usepackage[russian]{babel}
\usepackage{amsmath,amssymb,mathtools}
\usepackage{helvet}
\renewcommand{\familydefault}{\sfdefault}
\usepackage{icomma}
\usepackage[left=2cm,right=2cm,top=2cm,bottom=2.2cm,headheight=15pt]{geometry}
\usepackage[final]{microtype}
\usepackage{tikz}
\usetikzlibrary{arrows.meta,positioning}
\usepackage{tabularx,booktabs}
\usepackage{enumitem}
\usepackage{hyperref}
\usepackage{parskip}
\usepackage{fancyhdr}
\usepackage{setspace}
\usepackage{xcolor}
\usepackage{titlesec}

\definecolor{accent}{HTML}{008080}
\definecolor{darkaccent}{HTML}{004D4D}
\definecolor{textgray}{HTML}{333333}
\definecolor{softgray}{HTML}{E9EEEE}

\hypersetup{colorlinks=true,linkcolor=darkaccent,urlcolor=darkaccent}
\color{textgray}
\onehalfspacing
\setlength{\parindent}{0pt}
\setlength{\parskip}{0.55em}
\setlength{\emergencystretch}{2em}
\raggedbottom

\titleformat{\section}{\Large\bfseries\color{darkaccent}}{\thesection.}{0.6em}{}
\titleformat{\subsection}{\large\bfseries\color{accent}}{\thesubsection.}{0.5em}{}
\titlespacing*{\section}{0pt}{1.4em}{0.6em}
\titlespacing*{\subsection}{0pt}{1.1em}{0.4em}
\setlist[itemize]{leftmargin=1.7em,itemsep=0.25em,topsep=0.3em}
\setlist[enumerate]{leftmargin=2em,itemsep=0.7em,topsep=0.4em}

\pagestyle{fancy}
\fancyhf{}
\fancyhead[L]{\small\color{darkaccent}Математические модели}
\fancyhead[R]{\small\color{textgray}Кирилл Зиновьев}
\fancyfoot[C]{\small\thepage}
\renewcommand{\headrulewidth}{0.4pt}
\renewcommand{\headrule}{\hbox to\headwidth{\color{accent}\leaders\hrule height \headrulewidth\hfill}}

\newcommand{\important}[1]{\par\medskip{\color{darkaccent}\bfseries #1}\par\smallskip}

\begin{document}

\begin{titlepage}
\thispagestyle{empty}
\centering
\vspace*{2.8cm}
{\color{accent}\rule{0.72\linewidth}{1.1pt}\par}
\vspace{1.3cm}
{\fontsize{23}{28}\selectfont\bfseries Чек-лист\par}
\vspace{0.7cm}
{\fontsize{28}{34}\selectfont\bfseries\color{darkaccent}Составление математических моделей\par}
\vspace{0.45cm}
{\Large Дроби, проценты и буквенные выражения\par}
\vspace{1.5cm}
{\color{accent}\rule{0.42\linewidth}{0.8pt}\par}
\vfill
{\large Лёвин Артём Александрович\par}
\vspace{0.35cm}
{\large эксклюзивно для Кирилла Зиновьева\par}
\vspace{0.9cm}
{\large 7 класс\par}
\vspace{0.5cm}
{\large 22 августа 2026 года\par}
\vspace{2.1cm}
\begin{tikzpicture}[remember picture,overlay]
  \draw[color=accent,line width=1.4pt]
  ([xshift=1.5cm,yshift=1.8cm]current page.south west)
  .. controls ([xshift=5.3cm,yshift=3.05cm]current page.south west)
  and ([xshift=-5.7cm,yshift=0.8cm]current page.south east)
  .. ([xshift=-1.5cm,yshift=1.9cm]current page.south east);
\end{tikzpicture}
\end{titlepage}

\section{Что такое математическая модель}

\textbf{Математическая модель} --- это описание реальной ситуации с помощью чисел, букв, формул, уравнений, таблиц или схем.

Одна модель может описывать множество задач с разными сюжетами. Например, подписчиков социальной сети, учеников класса, стоимость товара и расход электроэнергии можно исследовать с помощью одинаковых формул.

\important{Главный принцип}

Буквы можно выбирать свободно. Каждая буква должна иметь точно определённый смысл, который сохраняется на протяжении всего решения.

\begin{center}
\begin{tikzpicture}[
  step/.style={draw=accent,rounded corners=2pt,minimum width=3cm,minimum height=1.25cm,align=center,fill=softgray,font=\small},
  arrow/.style={-{Stealth[length=2.5mm]},color=darkaccent,line width=0.8pt},
  node distance=0.65cm]
\node[step] (text) {Текст\\задачи};
\node[step,right=of text] (letters) {Величины\\и буквы};
\node[step,right=of letters] (formula) {Связи\\и формула};
\node[step,right=of formula] (check) {Вычисление\\и проверка};
\draw[arrow] (text) -- (letters);
\draw[arrow] (letters) -- (formula);
\draw[arrow] (formula) -- (check);
\end{tikzpicture}
\end{center}

\subsection{Алгоритм составления модели}

\begin{enumerate}
  \item Определите, какие величины известны и что требуется найти.
  \item Обозначьте величины буквами и запишите смысл каждой буквы.
  \item Установите связи между величинами.
  \item Составьте формулу или уравнение.
  \item Подставьте числовые данные, если они даны.
  \item Проверьте единицы измерения и разумность результата.
\end{enumerate}

\section{Нахождение дробной части числа}

Пусть (n) --- всё количество, а \(\dfrac{a}{b}\) --- выбранная доля, где \(b>0\) и \(0\le a\le b\). Тогда

\[
\boxed{P=\frac{a}{b}\,n},
\qquad
\boxed{R=n-P=\frac{b-a}{b}\,n}.
\]

\subsection*{Пример 1}

У Маши \(36\) подписчиков. Из них \(\dfrac{5}{6}\) смотрят научно-популярные видео:

\[
\frac{5}{6}\cdot36=30,
\qquad
36-30=6.
\]

\textbf{Ответ:} \(30\) подписчиков смотрят научно-популярные видео, \(6\) --- развлекательные ролики.

\section{Нахождение целого по известной части}

Пусть \(m\) --- известная дробная часть, \(\dfrac{a}{b}\) --- доля целого, \(N\) --- искомое целое, где \(a>0\), \(b>0\). Тогда

\[
\frac{a}{b}N=m,
\qquad
\boxed{N=m:\frac{a}{b}=m\cdot\frac{b}{a}}.
\]

\subsection*{Пример 2}

На олимпиаду пришли \(30\) семиклассников. Это составило \(\dfrac{3}{5}\) всех учеников параллели:

\[
N=30:\frac{3}{5}=30\cdot\frac{5}{3}=50.
\]

\textbf{Ответ:} \(50\) учеников.

\section{Математические модели для процентов}

\[
p\%=\frac{p}{100}.
\]

\subsection{Процент от числа}

\[
\boxed{\frac{p}{100}A}.
\]

\subsection{Уменьшение и увеличение}

\[
\boxed{A_1=A\left(1-\frac{p}{100}\right)},
\qquad
\boxed{A_1=A\left(1+\frac{q}{100}\right)}.
\]

Увеличение на \(400\%\) означает:

\[
A_1=A\left(1+\frac{400}{100}\right)=5A.
\]

Если величина \(A\) сначала уменьшилась на \(p\%\), затем увеличилась на \(q\%\) относительно полученного значения, то

\[
\boxed{A_2=A\left(1-\frac{p}{100}\right)\left(1+\frac{q}{100}\right)}.
\]

\subsection*{Пример 3}

В июле семья израсходовала \(300\) кВт\(\cdot\)ч электроэнергии. В августе расход уменьшился на \(70\%\), в сентябре увеличился на \(400\%\) по сравнению с августом.

\[
A_1=300\left(1-\frac{70}{100}\right)=90\ \text{кВт}\cdot\text{ч},
\]

\[
A_2=90\left(1+\frac{400}{100}\right)=450\ \text{кВт}\cdot\text{ч}.
\]

\textbf{Ответ:} июль --- \(300\), август --- \(90\), сентябрь --- \(450\) кВт\(\cdot\)ч.

\section{Модели с несколькими слагаемыми}

Если общее количество складывается из нескольких частей, модель строится в виде суммы.

\subsection*{Пример 4. Урожай с двух участков}

Площадь первого участка равна \(a\) га, второго --- \(b\) га. Урожайность составляет соответственно \(32\) и \(40\) ц/га. Общий урожай:

\[
\boxed{S=32a+40b}.
\]

При \(a=120\) и \(b=80\):

\[
S=32\cdot120+40\cdot80=7040\ \text{ц}.
\]

\section{Как проверить составленную модель}

\begin{itemize}[label=\(\square\)]
  \item Каждая буква получила точное объяснение.
  \item Одинаковые величины измеряются в согласованных единицах.
  \item Знаменатели отличны от нуля.
  \item Доля целого находится умножением целого на дробь.
  \item Целое по известной дробной части находится делением части на дробь.
  \item При уменьшении используется множитель \(1-\dfrac{p}{100}\).
  \item При увеличении используется множитель \(1+\dfrac{p}{100}\).
  \item Каждый следующий процент вычисляется от значения, указанного в условии.
  \item Полученный результат соответствует смыслу задачи.
\end{itemize}

\section{Типичные ошибки}

\begin{itemize}
  \item \textbf{Подмена целого его частью.} В задаче следует отдельно обозначить всё количество и выбранную долю.
  \item \textbf{Остановка после вычисления величины изменения.} Число \(\dfrac{p}{100}A\) показывает, насколько изменилась величина. Итоговое значение требует сложения или вычитания.
  \item \textbf{Вычисление всех процентов от исходного числа.} При последовательных изменениях база каждого процента определяется условием.
  \item \textbf{Неверное понимание выражения «увеличилось на \(p\%\)».} Итоговый множитель равен \(1+\dfrac{p}{100}\).
  \item \textbf{Потеря единиц измерения.} Например, \(\dfrac{\text{ц}}{\text{га}}\cdot\text{га}=\text{ц}\).
\end{itemize}

\clearpage
\section*{Тренировочный блок}

\textbf{Указание.} Для каждой задачи запишите математическую модель, затем выполните вычисления.

\subsection*{Средний уровень}

\begin{enumerate}
  \item В группе \(54\) ученика. Из них \(\dfrac{7}{9}\) посещают математический кружок. Сколько учеников посещают кружок? Сколько учеников его пока не посещают?
  \item На соревнования пришли \(28\) учеников, что составило \(\dfrac{4}{7}\) всех участников школьной команды. Сколько учеников входит в команду?
  \item В библиотеку поступило \(80\) книг. Детективы составили \(\dfrac{3}{8}\) всех книг, остальные книги были научно-популярными. Найдите количество книг каждого вида.
\end{enumerate}

\subsection*{Уровень выше среднего}

\begin{enumerate}
  \setcounter{enumi}{3}
  \item Цена прибора составляла \(2400\) рублей. Сначала её уменьшили на \(15\%\), затем новую цену увеличили на \(10\%\). Найдите окончательную цену.
  \item Предприятие расходовало \(500\) м\(^3\) воды в месяц. После модернизации расход уменьшился на \(24\%\), затем из-за увеличения производства вырос на \(25\%\) относительно нового значения. Найдите конечный расход.
  \item Из \(N\) участников доля \(\dfrac{a}{b}\) выбрала первый маршрут. Запишите формулы для количества участников первого и второго маршрутов. Укажите ограничения на \(a\) и \(b\).
  \item Площадь первого поля равна \(90\) га, второго --- \(60\) га. Урожайность составила соответственно \(28\) и \(35\) ц/га. Составьте буквенную модель общего урожая и вычислите его.
  \item В магазин привезли \(x\) упаковок по \(12\) тетрадей и \(y\) упаковок по \(18\) тетрадей. Составьте формулу общего количества тетрадей. Вычислите значение при \(x=15\) и \(y=10\).
  \item После уменьшения числа на \(p\%\) полученный результат увеличили на \(25\%\). В итоге вернулись к исходному числу. Найдите \(p\).
\end{enumerate}

\subsection*{Сложный уровень}

\begin{enumerate}
  \setcounter{enumi}{9}
  \item В параллели \(\dfrac{2}{5}\) учеников посещают спортивные секции. Из оставшихся учеников \(30\%\) занимаются музыкой. Остальные \(42\) ученика пока не выбрали кружок. Составьте математическую модель и найдите общее количество учеников, число спортсменов и число музыкантов.
\end{enumerate}

\clearpage
\section*{Ответы к тренировочным упражнениям}

\renewcommand{\arraystretch}{1.45}
\begin{table}[ht]
\centering
\caption{Краткие ответы}
\begin{tabularx}{\linewidth}{@{}cX@{}}
\toprule
\textbf{№} & \textbf{Ответ} \\
\midrule
1 & \(42\) ученика посещают кружок; \(12\) учеников пока его не посещают. \\
2 & \(49\) учеников. \\
3 & \(30\) детективов и \(50\) научно-популярных книг. \\
4 & \(2244\) руб. \\
5 & \(475\) м\(^3\). \\
6 & \(\dfrac{a}{b}N\) и \(\dfrac{b-a}{b}N\); \(b>0\), \(0\le a\le b\). \\
7 & \(4620\) ц. \\
8 & \(360\) тетрадей. \\
9 & \(20\%\). \\
10 & Всего \(100\) учеников; спортсменов \(40\), музыкантов \(18\), кружок пока не выбрали \(42\). \\
\bottomrule
\end{tabularx}
\end{table}

\end{document}
